% !TEX program = lualatex
\documentclass[aspectratio=43,professionalfonts,handout]{beamer}
\usepackage{beamer_template}

\newcommand{\LectureNum}{12}
\newcommand{\LectureTitle}{フーリエ変換と積分定理}
\newcommand{\TextPages}{p.\,91-93}
\newcommand{\CourseName}{応用数学}
\renewcommand{\TermName}{前期}

\begin{document}

\begin{frame}{\CourseName\quad 第\LectureNum 回「\LectureTitle」}
  {\small \TermName\quad \TeacherName\quad ／\quad 教科書 \TextPages}

  \textbf{目標}：フーリエ変換の定義と積分定理を理解し、基本的なフーリエ変換を求めることができる

\end{frame}

\begin{frame}{フーリエ変換の定義}
  \begin{block}{}
  \begin{itemize}
    \item 非周期関数への拡張（周期 $ \to \infty$ ）
    \item $F(u) = \int_{-\infty}^{\infty} f(x)e^{-iux} dx$
    \item ラプラス変換との比較
  \end{itemize}
  \end{block}
\end{frame}

\begin{frame}{フーリエ積分定理・反転公式}
  \begin{block}{}
  \begin{itemize}
    \item 積分定理：$f(x) \approx (1/2\pi)\int F(u)e^{iux}du$
    \item 反転公式 $f = \mathcal{F}^{-1}[F]$
    \item 連続点では等号成立
  \end{itemize}
  \end{block}
\end{frame}

\begin{frame}{例題1・2}
  \begin{block}{}
  \begin{itemize}
    \item 例題1：$f(x)=e^{-|x|}  \to  F(u)=2/(1+u^{2})$
    \item 例題2：矩形波のフーリエ変換
    \item 余弦・正弦変換（偶関数・奇関数）
  \end{itemize}
  \end{block}
\end{frame}

\begin{frame}{演習}
  \begin{exampleblock}{自分で解いてみよう}
  \begin{itemize}
    \item 問1〜4より選題
  \end{itemize}
  \end{exampleblock}
\end{frame}

\begin{frame}{まとめ}
  \begin{itemize}
    \item フーリエ変換の定義
    \item 積分定理・反転公式
    \item 次回：フーリエ変換の性質
  \end{itemize}
\end{frame}

\end{document}
